\documentclass[11pt]{article}

\usepackage{amsmath,amssymb}
\usepackage[margin=1in]{geometry}

\newcommand{\R}{\mathbb{R}}
\newcommand{\Om}{\Omega}
\newcommand{\dd}{\mathrm{d}}

\begin{document}

\section*{Examples of Weak Derivatives}

Recall that a function $w\in L^1_{\mathrm{loc}}(\Om)$ is the
\emph{weak derivative} of $u\in L^1_{\mathrm{loc}}(\Om)$ if
\begin{equation}
    \int_{\Om} w\varphi\,\dd x
    =
    -\int_{\Om}u\frac{\dd\varphi}{\dd x}\,\dd x
    \qquad
    \forall \varphi\in C_0^\infty(\Om).
    \label{eq:weak-derivative}
\end{equation}
We then write
\[
w=\frac{\dd u}{\dd x},
\]
where the derivative is understood in the weak sense.

\subsection*{Example 1: The absolute-value function}

Let
\[
\Om=(-1,1),
\qquad
u(x)=|x|.
\]
The classical derivative of $u$ does not exist at $x=0$. Define
\[
w(x)=
\begin{cases}
-1, & -1<x<0,\\[2mm]
 1, & 0<x<1.
\end{cases}
\]
We claim that $w$ is the weak derivative of $u$.

To prove the claim, let $\varphi\in C_0^\infty(-1,1)$. We must show that
\[
\int_{-1}^{1}w\varphi\,\dd x
=
-\int_{-1}^{1}u\frac{\dd\varphi}{\dd x}\,\dd x.
\]

For the left-hand side,
\begin{align}
\int_{-1}^{1}w\varphi\,\dd x
&=
\int_{-1}^{0}(-1)\varphi\,\dd x
+
\int_{0}^{1}\varphi\,\dd x
\nonumber\\
&=
-\int_{-1}^{0}\varphi\,\dd x
+
\int_{0}^{1}\varphi\,\dd x.
\label{eq:absolute-lhs}
\end{align}

For the right-hand side, since
\[
|x|=
\begin{cases}
-x, & -1<x<0,\\
 x, & 0<x<1,
\end{cases}
\]
we have
\begin{align}
-\int_{-1}^{1}|x|\frac{\dd\varphi}{\dd x}\,\dd x
&=
-\int_{-1}^{0}(-x)\frac{\dd\varphi}{\dd x}\,\dd x
-\int_{0}^{1}x\frac{\dd\varphi}{\dd x}\,\dd x
\nonumber\\
&=
\int_{-1}^{0}x\frac{\dd\varphi}{\dd x}\,\dd x
-\int_{0}^{1}x\frac{\dd\varphi}{\dd x}\,\dd x.
\label{eq:absolute-rhs-1}
\end{align}

Using the product rule,
\[
x\frac{\dd\varphi}{\dd x}
=
\frac{\dd}{\dd x}(x\varphi)-\varphi,
\]
so
\begin{align}
-\int_{-1}^{1}|x|\frac{\dd\varphi}{\dd x}\,\dd x
&=
\int_{-1}^{0}
\left[
\frac{\dd}{\dd x}(x\varphi)-\varphi
\right]\dd x
-
\int_{0}^{1}
\left[
\frac{\dd}{\dd x}(x\varphi)-\varphi
\right]\dd x
\nonumber\\
&=
\left[x\varphi(x)\right]_{-1}^{0}
-\int_{-1}^{0}\varphi\,\dd x
-\left[x\varphi(x)\right]_{0}^{1}
+\int_{0}^{1}\varphi\,\dd x.
\label{eq:absolute-rhs-2}
\end{align}

Since $\varphi\in C_0^\infty(-1,1)$,
\[
\varphi(-1)=\varphi(1)=0.
\]
The terms involving $\varphi(0)$ also vanish because they are
multiplied by $x=0$. Therefore,
\[
\left[x\varphi(x)\right]_{-1}^{0}
=
\left[x\varphi(x)\right]_{0}^{1}
=
0.
\]
Hence,
\[
-\int_{-1}^{1}|x|\frac{\dd\varphi}{\dd x}\,\dd x
=
-\int_{-1}^{0}\varphi\,\dd x
+
\int_{0}^{1}\varphi\,\dd x.
\]
Comparing with \eqref{eq:absolute-lhs},
\[
\int_{-1}^{1}w\varphi\,\dd x
=
-\int_{-1}^{1}u\frac{\dd\varphi}{\dd x}\,\dd x.
\]
Thus, $w$ is the weak derivative of $u$, and
\[
\boxed{
\frac{\dd u}{\dd x}
=
\begin{cases}
-1, & -1<x<0,\\[2mm]
 1, & 0<x<1,
\end{cases}
\quad\text{a.e. in }(-1,1).
}
\]

The value assigned to the weak derivative at the single point $x=0$
is irrelevant.

\subsection*{Example 2: Strain jump at a material interface}

Weak derivatives arise naturally in solid mechanics. Consider a
one-dimensional bar occupying
\[
\Om=(0,L)
\]
and composed of two linear elastic materials joined at $x=a$, where
$0<a<L$. The Young's modulus is
\[
E(x)=
\begin{cases}
E_1, & 0<x<a,\\[2mm]
E_2, & a<x<L,
\end{cases}
\qquad
E_1\neq E_2.
\]

Assume that the bar has constant cross-sectional area $A$, is fixed
at $x=0$, and is subjected to an axial end load $P$ at $x=L$.
The axial stress is constant throughout the bar:
\[
\sigma_{xx}=\frac{P}{A}.
\]
From Hooke's law,
\[
\sigma_{xx}=E\varepsilon_{xx},
\]
the strain is
\[
\varepsilon_{xx}(x)=
\begin{cases}
\dfrac{P}{AE_1}, & 0<x<a,\\[3mm]
\dfrac{P}{AE_2}, & a<x<L.
\end{cases}
\]
Thus,
\[
\varepsilon_{xx}(a^-)
=
\frac{P}{AE_1}
\neq
\frac{P}{AE_2}
=
\varepsilon_{xx}(a^+),
\]
so the strain has a jump at the material interface.

Within each material,
\[
\varepsilon_{xx}=\frac{\dd u}{\dd x}.
\]
Using $u(0)=0$ and displacement continuity at $x=a$, the displacement
field is
\[
u(x)=
\begin{cases}
\dfrac{P}{AE_1}x,
&0\leq x\leq a,\\[4mm]
\dfrac{P}{AE_1}a
+\dfrac{P}{AE_2}(x-a),
&a<x\leq L.
\end{cases}
\]

The displacement is continuous at the material interface:
\[
u(a^-)=u(a^+)=\frac{Pa}{AE_1}.
\]
However, its one-sided classical derivatives are
\[
\left.\frac{\dd u}{\dd x}\right|_{x=a^-}
=
\frac{P}{AE_1},
\qquad
\left.\frac{\dd u}{\dd x}\right|_{x=a^+}
=
\frac{P}{AE_2}.
\]
Since $E_1\neq E_2$, these two values are different. Therefore, the
classical derivative of $u$ does not exist at $x=a$.

Define
\[
w(x)=
\begin{cases}
\dfrac{P}{AE_1}, & 0<x<a,\\[3mm]
\dfrac{P}{AE_2}, & a<x<L.
\end{cases}
\]
We claim that $w$ is the weak derivative of $u$. Physically,
\[
w=\varepsilon_{xx}.
\]

To prove the claim, let $\varphi\in C_0^\infty(0,L)$. We must show that
\[
\int_0^L w\varphi\,\dd x
=
-\int_0^L u\frac{\dd\varphi}{\dd x}\,\dd x.
\]

Split the right-hand side at the material interface:
\begin{align}
-\int_0^L u\frac{\dd\varphi}{\dd x}\,\dd x
&=
-\int_0^a u\frac{\dd\varphi}{\dd x}\,\dd x
-\int_a^L u\frac{\dd\varphi}{\dd x}\,\dd x.
\label{eq:bar-split}
\end{align}

On each material subdomain, $u$ is classically differentiable.
Using integration by parts,
\begin{align}
-\int_0^a u\frac{\dd\varphi}{\dd x}\,\dd x
&=
-\left[u\varphi\right]_0^a
+
\int_0^a \frac{\dd u}{\dd x}\varphi\,\dd x
\nonumber\\
&=
-u(a^-)\varphi(a)
+u(0)\varphi(0)
+
\int_0^a \frac{P}{AE_1}\varphi\,\dd x,
\label{eq:bar-left}
\end{align}
and
\begin{align}
-\int_a^L u\frac{\dd\varphi}{\dd x}\,\dd x
&=
-\left[u\varphi\right]_a^L
+
\int_a^L \frac{\dd u}{\dd x}\varphi\,\dd x
\nonumber\\
&=
-u(L)\varphi(L)
+u(a^+)\varphi(a)
+
\int_a^L \frac{P}{AE_2}\varphi\,\dd x.
\label{eq:bar-right}
\end{align}

Since $\varphi\in C_0^\infty(0,L)$,
\[
\varphi(0)=\varphi(L)=0.
\]
Moreover, displacement continuity at the material interface gives
\[
u(a^-)=u(a^+).
\]
Therefore, the interface terms cancel:
\[
-u(a^-)\varphi(a)+u(a^+)\varphi(a)=0.
\]

Adding \eqref{eq:bar-left} and \eqref{eq:bar-right} gives
\begin{align}
-\int_0^L u\frac{\dd\varphi}{\dd x}\,\dd x
&=
\int_0^a \frac{P}{AE_1}\varphi\,\dd x
+
\int_a^L \frac{P}{AE_2}\varphi\,\dd x
\nonumber\\
&=
\int_0^L w\varphi\,\dd x.
\end{align}
Hence, $w$ is the weak derivative of $u$, and
\[
\boxed{
\frac{\dd u}{\dd x}
=
\varepsilon_{xx}
=
\begin{cases}
\dfrac{P}{AE_1}, & 0<x<a,\\[3mm]
\dfrac{P}{AE_2}, & a<x<L,
\end{cases}
\quad\text{a.e. in }(0,L).
}
\]

The value assigned to the weak derivative at the single point $x=a$
is irrelevant.

Thus, the displacement $u$ is continuous but is not classically
differentiable at the material interface. Nevertheless, it possesses
a weak derivative, and this weak derivative is precisely the
piecewise-defined strain $\varepsilon_{xx}$. This example illustrates
why weak derivatives arise naturally in mechanics when material
properties are discontinuous.

\end{document}